\section{Step 4: triage -- from a failure to a bug report}
\label{sec:triage}

A campaign produces failure directories (one per new signature). Most
candidates are not reportable as they are: test cases are long, some are
flaky, some are our own false positives. \texttt{python -m pgfuzz triage
<config> <failure dir>} (\file{src/pgfuzz/triage/triage.py}) turns a candidate
into either a verdict \emph{flaky} or a reduced reproducer with a draft report.

\begin{center}
\begin{tikzpicture}[font=\small, node distance=4mm,
  b/.style={draw,rounded corners=2pt,fill=blue!4,align=center,minimum height=9mm,text width=21mm}]
\node[b] (a) {reproduce $\times 3$};
\node[b,right=of a] (b) {tree reduction};
\node[b,right=of b] (c) {statement ddmin};
\node[b,right=of c] (d) {token HDD (pairs)};
\node[b,right=of d] (e) {observe + localize};
\node[b,right=of e] (f) {report, upstream draft};
\foreach \x/\y in {a/b,b/c,c/d,d/e,e/f} \draw[-{Stealth}] (\x)--(\y);
\end{tikzpicture}
\end{center}

\subsection{Reproduce}

The candidate is executed three times by the \emph{same executor and oracle}
as the campaign, configured for its oracle: only the failing plan variant, or
only its metamorphic check (with confirmation on). \texttt{same\_failure}
decides whether a re-run shows \emph{the same} failure, not just any failure:
\begin{itemize}
  \item assertion: same file and same expression;
  \item wrong result: same variant / oracle kind;
  \item internal error, corruption, panic: same normalised message;
  \item crash without assertion: any crash of the same statement.
\end{itemize}
0/3 $\Rightarrow$ verdict \emph{flaky}, nothing else is done (2 of our 18
candidates). This criterion is also what every reduction step uses.

\subsection{Reduction, three levels}
\label{sec:reduce}

\paragraph{1. Tree reduction} (\file{src/pgfuzz/triage/tree_reduce.py}) -- the
fourth job of the derivation tree. The failure directory stores the trees of
the test case (JSON, stable alternative ordinals). Moves, biggest expected gain
first:
\begin{enumerate}
  \item drop a whole statement;
  \item drop one element of a repetition (\texttt{\#rep});
  \item replace a subtree by the \emph{smallest derivation of the same rule and
    type}: the node is replaced by a \texttt{MINIMAL} marker and the engine
    derives it with its minimal-depth alternatives.
\end{enumerate}
After each move the generator \emph{replays} the trees (re-binding names,
aliases and types), so almost every candidate is valid SQL: the server's time
goes to the question that matters (does the failure remain?), not to rejecting
broken SQL. A candidate is only tried if it is smaller in tokens; after each
success the search restarts from the biggest moves.

\paragraph{2. Statement ddmin} (\texttt{reduce\_statements}): remove setup
statements that are not needed (Zeller \& Hildebrandt's delta debugging),
always keeping the culprit.

\paragraph{3. Token HDD} (\texttt{TokenReducer}, \file{reduce.py}): hierarchical
delta debugging over the \emph{parenthesis tree} of the culprit statement
(Misherghi \& Su). Coarse passes first: remove a clause (\texttt{WHERE},
\texttt{ORDER BY}, \dots), remove one list element, unwrap a parenthesised
sub-expression or replace it by a small literal (\texttt{NULL}, \texttt{1},
\texttt{TRUE}, \texttt{FALSE}, \texttt{0}, \texttt{''}); then single tokens.
Results are cached by text; budget 500 tests or 15 minutes.

\paragraph{Pairs.} For a metamorphic candidate (query $Q$ and rewrite $Q'$)
both sides must stay equivalent while being reduced. \texttt{split\_pair}
computes the common prefix and suffix of the two token lists and leaves a
\texttt{HOLE} where they differ; the reducer edits the shared skeleton, and both
sides are materialised by filling the hole with each side's middle. So every
edit is applied to both sides at once and the hole is never removed.

\begin{keybox}[Measured: tree reduction vs token reduction (BUG-01)]
A real grammar test case failing BUG-01's assertion: 3 statements, 590 tokens.
\begin{center}
\begin{tabular}{lrrr}
\toprule
Reduction & Final culprit & Server runs & Time \\
\midrule
statements + tokens only (\texttt{--no-tree}) & 116 tokens & 406 & 33.7 s \\
tree first, then statements + tokens & \textbf{79 tokens} & 555 & 39.6 s \\
\bottomrule
\end{tabular}
\end{center}
The tree stage alone: 590 $\rightarrow$ 136 tokens in 150 runs (196 replays).
Token reduction gets stuck in local minima (removing any single token breaks
the SQL); the tree can jump to a \emph{different, smaller} valid derivation.
Evidence: \file{docs/evidence/tree_reduction/}. BUG-02's first candidate:
4 statements / 107 tokens $\rightarrow$ 1 statement / 36 tokens in 307 runs
(token level, before tree reduction existed).
\end{keybox}

\subsection{Observe and localize}
\label{sec:localize}

After reduction, triage runs the queries to show (both sides for a pair, or
the query under default and failing settings) and records up to 30 result rows
on the test server and, if configured, on the 18.6 reference.

Then \file{src/pgfuzz/triage/localize.py} combines the two oracle families
(plan section~0.3 ``plan oracle $\times$ rewrite oracle''): it re-runs the
reduced queries under \emph{every} plan variant and reports, per variant, whether
the answer changes and whether the two sides agree. Verdicts:
\begin{description}
  \item[planner] the disagreement disappears when some settings are switched
    off: those code paths are the first suspects. A table
    (\texttt{SUSPECTS}) maps each setting to source files, e.g.
    \texttt{enable\_hashjoin} $\rightarrow$ \file{executor/nodeHashjoin.c,
    optimizer/path/joinpath.c}; \texttt{enable\_eager\_aggregate}
    $\rightarrow$ \file{optimizer/path/allpaths.c, optimizer/util/relnode.c}.
  \item[not toggled by planner settings] no setting removes it: suspect an
    always-applied rewrite, constant folding, a planner support function or
    expression evaluation. \textbf{BUG-02 is localized this way}, and its true
    cause is exactly a planner support function (\texttt{rangetypes.c}).
  \item[sides agree by default] the reduction went wrong (a warning).
\end{description}

\subsection{Outputs}

Per candidate, in \file{results/triage/<signature>/}:
\begin{itemize}
  \item \file{reduced.sql}: setup statements, then the culprit; for plan
    variants the \texttt{SET} commands are written out; for metamorphic
    candidates both queries, so the file is a stand-alone reproducer for
    \texttt{psql};
  \item \file{triage.json}: reproduction counts, sizes before/after, runs,
    seconds, observed rows, localization, server versions;
  \item \file{report.md}: a readable report;
  \item \file{upstream_report.txt}: a draft for the \texttt{pgsql-bugs} form,
    in the shape \pg developers expect (version and build, what happens, the
    minimal script, actual output on master and on the release, the
    localization, the top 12 backtrace frames, and two placeholders we must
    fill by hand: \emph{expected behaviour with a justification from the SQL
    standard or the documentation}, and the bisected commit).
\end{itemize}
The tool never writes the ``expected'' part: deciding that something is a bug
is a human judgement.

\paragraph{Bisection.} \file{scripts/bisect_pg.sh GOOD BAD reduced.sql MODE}
runs \texttt{git bisect} on a full-history \pg clone, building each step into a
scratch prefix. Modes: \texttt{crash} (backend crash or assertion),
\texttt{error} (psql reports an ERROR), \texttt{output} (differs from an expected
output), \texttt{differ} (the last two queries of \file{reduced.sql} return
different rows -- for metamorphic and plan-variant candidates). Commits that do
not build are skipped (exit 125).

\section{The precision ledger (RQ2 and claim 2)}
\label{sec:ledger}

Every triaged candidate gets an entry in \file{triage/verdicts.yaml}: signature,
oracle, verdict (\emph{real, duplicate, false\_positive, not\_a\_bug, flaky}),
a note, \texttt{fixed: true} if the false-positive class has been fixed in the
tool, and the upstream link. \texttt{python -m pgfuzz precision} computes
(\file{src/pgfuzz/triage/precision.py}):
\[
\text{precision} = \frac{\text{real}}{\text{real} + \text{false positive} + \text{not a bug}},
\qquad
\text{after fixes} = \frac{\text{real}}{\text{decided} - \text{fixed}}.
\]
Duplicates count towards neither side; flaky entries are reported separately.

\paragraph{The ledger today (18 candidates).}
\begin{center}
\small
\begin{tabular}{lrrrrrr}
\toprule
Oracle & Candidates & Real & False pos. & Duplicates & Precision & After fixes \\
\midrule
assertion & 2 & 1 & 1 & 0 & 0.50 & 1.00 \\
crash & 1 & 0 & 1 & 0 & 0.00 & -- \\
cross-version & 3 & 0 & 2 & 0 & 0.00 & -- \\
internal\_error & 2 & 1 & 0 & 0 & 1.00 & 1.00 \\
metamorphic:expr & 1 & 1 & 0 & 0 & 1.00 & 1.00 \\
metamorphic:tlp\_having & 4 & 0 & 4 & 0 & 0.00 & -- \\
plan-variant & 5 & 0 & 5 & 0 & 0.00 & -- \\
\midrule
all & 18 & 3 & 13 & 0 & 0.19 & 1.00 \\
\bottomrule
\end{tabular}

\end{center}

\begin{center}
\scriptsize
\begin{longtable}{p{0.16\linewidth}p{0.17\linewidth}p{0.12\linewidth}p{0.47\linewidth}}
\toprule
Signature & Oracle & Verdict & Note \\
\midrule
\endhead
\rows{ledger_rows.tex}
\bottomrule
\end{longtable}
\end{center}

\paragraph{How to read it (be ready for this question).}
\begin{itemize}
  \item The 3 real ones: BUG-01 (assertion), BUG-02 (metamorphic expression
    rewrite), and the re-injected benchmark bug \texttt{rb-dfcf6fa7} (internal
    error), which counts as real because it is a real historical \pg bug, put
    back on purpose.
  \item 13 false positives / not-a-bug come from \textbf{6 classes}, all fixed
    in the tool with a test: (1) \texttt{ARRAY(SELECT)} without
    \texttt{ORDER BY} (5 plan-variant candidates); (2) clock-dependent seed
    data across versions; (3) TLP-HAVING without \texttt{GROUP BY} (4);
    (4) \texttt{DISTINCT} over numeric scale; (5) the shared-server PoC replay
    (2, assertion + crash); (6) the intended \texttt{corr} change (not a bug).
  \item Raw precision 0.19 is low \emph{on purpose}: we log every false
    positive we ever produced, including during development. ``After fixes''
    1.00 means every remaining decided report of a class we have not fixed was
    real. The paper reports both and the per-oracle counts of the long
    campaigns, which will be the honest number.
  \item The 798 new-kind checks, the 532 RLS checks, the 58 CERT checks and the
    1{,}991 metamorphic checks of the full-system smoke produced 0 new false
    positives.
\end{itemize}
